% GENERATED FILE. Edit canonical lessons, then run npm run generate:books.
% canonical-source-hash: fnv1a64:2bf9c3973b0c8316
% canonical-lessons: SA-C243-balavan, SA-C243-rugnah, SA-C243-svasthah, SA-C243-sukhi, SA-C243-duhkhi

\chapter{Mighty, Sick, Healthy, Happy, Unhappy}
\label{ch:sa-mighty-sick-healthy-happy-unhappy}
\hlchaptermodality{243}

\begin{chapteropening}
\textbf{By the end of this chapter:} \emph{I can say mighty, sick, healthy, happy and unhappy.}

Complete the last lesson of chapter 243: I can say mighty, sick, healthy, happy and unhappy.
\end{chapteropening}

\section[\texorpdfstring{\sk{बलवान्} (balavān)}{बलवान् (balavān)}]{\sk{बलवान्} (balavān) — mighty, strong}

\label{lesson:SA-C243-balavan}



\begin{quote}
Before the new one: say the Sanskrit for easy to do, then the Sanskrit for hard to do.
\end{quote}

\subsection*{You'll want to know: \sk{बलवान्}}
\textbf{\sk{बलवान्}} — \emph{balavān} — \textquotedblleft{}mighty, strong\textquotedblright{}.

\begin{grammarlens}[title={What you've built}]
One more word for this chapter.
\end{grammarlens}

\subsection*{Your turn}
\begin{itemize}
\raggedright
  \item \emph{Say it:} \emph{balavān}
  \item \emph{Say it:} \emph{balavān}, once more
  \item \emph{Say it:} say \emph{balavān}
  \item \emph{Recall:} say the Sanskrit for easy to do, then the Sanskrit for hard to do, then say \emph{balavān} again
\end{itemize}

\begin{tcolorbox}[breakable,colback=teal!4,colframe=teal!35!black,title={Before you move on}]
What does \sk{बलवान्} mean? (\textquotedblleft{}Mighty, strong\textquotedblright{}.) Say it once more.
\end{tcolorbox}
\section[\texorpdfstring{\sk{रुग्णः} (rugṇaḥ)}{रुग्णः (rugṇaḥ)}]{\sk{रुग्णः} (rugṇaḥ) — sick}

\label{lesson:SA-C243-rugnah}



\begin{quote}
Before the new one: say the Sanskrit for hard to do, then the Sanskrit for mighty.
\end{quote}

\subsection*{You'll want to know: \sk{रुग्णः}}
\textbf{\sk{रुग्णः}} — \emph{rugṇaḥ} — \textquotedblleft{}sick\textquotedblright{}.

\begin{grammarlens}[title={What you've built}]
One more word for this chapter.
\end{grammarlens}

\subsection*{Your turn}
\begin{itemize}
\raggedright
  \item \emph{Say it:} \emph{rugṇaḥ}
  \item \emph{Say it:} \emph{rugṇaḥ}, once more
  \item \emph{Say it:} say \emph{rugṇaḥ}
  \item \emph{Recall:} say the Sanskrit for hard to do, then the Sanskrit for mighty, then say \emph{rugṇaḥ} again
\end{itemize}

\begin{tcolorbox}[breakable,colback=teal!4,colframe=teal!35!black,title={Before you move on}]
What does \sk{रुग्णः} mean? (\textquotedblleft{}Sick\textquotedblright{}.) Say it once more.
\end{tcolorbox}
\section[\texorpdfstring{\sk{स्वस्थः} (svasthaḥ)}{स्वस्थः (svasthaḥ)}]{\sk{स्वस्थः} (svasthaḥ) — healthy}

\label{lesson:SA-C243-svasthah}



\begin{quote}
Before the new one: say the Sanskrit for mighty, then the Sanskrit for sick.
\end{quote}

\subsection*{You'll want to know: \sk{स्वस्थः}}
\textbf{\sk{स्वस्थः}} — \emph{svasthaḥ} — \textquotedblleft{}healthy\textquotedblright{}.

\begin{grammarlens}[title={What you've built}]
One more word for this chapter.
\end{grammarlens}

\subsection*{Your turn}
\begin{itemize}
\raggedright
  \item \emph{Say it:} \emph{svasthaḥ}
  \item \emph{Say it:} \emph{svasthaḥ}, once more
  \item \emph{Say it:} say \emph{svasthaḥ}
  \item \emph{Recall:} say the Sanskrit for mighty, then the Sanskrit for sick, then say \emph{svasthaḥ} again
\end{itemize}

\begin{tcolorbox}[breakable,colback=teal!4,colframe=teal!35!black,title={Before you move on}]
What does \sk{स्वस्थः} mean? (\textquotedblleft{}Healthy\textquotedblright{}.) Say it once more.
\end{tcolorbox}
\section[\texorpdfstring{\sk{सुखी} (sukhī)}{सुखी (sukhī)}]{\sk{सुखी} (sukhī) — happy}

\label{lesson:SA-C243-sukhi}



\begin{quote}
Before the new one: say the Sanskrit for sick, then the Sanskrit for healthy.
\end{quote}

\subsection*{You'll want to know: \sk{सुखी}}
\textbf{\sk{सुखी}} — \emph{sukhī} — \textquotedblleft{}happy\textquotedblright{}.

\begin{grammarlens}[title={What you've built}]
One more word for this chapter.
\end{grammarlens}

\subsection*{Your turn}
\begin{itemize}
\raggedright
  \item \emph{Say it:} \emph{sukhī}
  \item \emph{Say it:} \emph{sukhī}, once more
  \item \emph{Say it:} say \emph{sukhī}
  \item \emph{Recall:} say the Sanskrit for sick, then the Sanskrit for healthy, then say \emph{sukhī} again
\end{itemize}

\begin{tcolorbox}[breakable,colback=teal!4,colframe=teal!35!black,title={Before you move on}]
What does \sk{सुखी} mean? (\textquotedblleft{}Happy\textquotedblright{}.) Say it once more.
\end{tcolorbox}
\section[\texorpdfstring{\sk{दुःखी} (duḥkhī)}{दुःखी (duḥkhī)}]{\sk{दुःखी} (duḥkhī) — unhappy}

\label{lesson:SA-C243-duhkhi}



\begin{quote}
Before the new one: say the Sanskrit for healthy, then the Sanskrit for happy.
\end{quote}

\subsection*{You'll want to know: \sk{दुःखी}}
\textbf{\sk{दुःखी}} — \emph{duḥkhī} — \textquotedblleft{}unhappy\textquotedblright{}.

\begin{grammarlens}[title={What you've built}]
One more word for this chapter.
\end{grammarlens}

\subsection*{Your turn}
\begin{itemize}
\raggedright
  \item \emph{Say it:} \emph{duḥkhī}
  \item \emph{Say it:} \emph{duḥkhī}, once more
  \item \emph{Say it:} say \emph{duḥkhī}
  \item \emph{Recall:} say the Sanskrit for healthy, then the Sanskrit for happy, then say \emph{duḥkhī} again
\end{itemize}

\begin{tcolorbox}[breakable,colback=teal!4,colframe=teal!35!black,title={Before you move on}]
What does \sk{दुःखी} mean? (\textquotedblleft{}Unhappy\textquotedblright{}.) Say it once more.
\end{tcolorbox}
